\section{Evaluation}\label{sec:evaluation}

\subsection{Evaluation Setup}

\paragraph{Hardware.} We implement \textsc{Terrae} in C++ and evaluate
it on on Ubuntu 22.04.5 LTS (Linux kernel 6.5.0-18-generic, x86\_64).
The evaluation machine is equipped with Intel Xeon Gold 5122 CPUs at 3. 
60 GHz, providing 32 logical  
cores in total, and 503 GiB of RAM.

\paragraph{Baseline.}
There are two candidates for baseline: (1) Sparrow~\cite{sparrow24}, and (2) ZKBoost~\cite{zkboost26}. The first is 
designed for random forests, and the second is designed for XGBoost,
a specific implementation of GBDT.
The two systems share a same design philosophy for
committing forest and proving training and inference.
We adopt ZKBoost as our baseline, as it is more closely aligned with our design and can be easily adapted to support GBDT.
To ensure a fair comparison, we must keep the core design of ZKBoost,
which include
\begin{enumerate}
    \item The approach of committing forest.
        In ZKBoost, each tree is padded with
        dummy nodes to the same depth
        and to a full binary tree.
        This makes the proof of routing more easy,
        since the structure of the forest is fixed and public.
    \item The approach of proving inference.
        In ZKBoost, the prover commits the paths
        the samples take
        from the roots to the leaves,
        validates that the paths are consistent with the committed forest,
        and then proves the correctness of the leaf values.
        This follows the style
        of \cite{zkdt20} for decision tree inference.
    \item The approach of proving training.
        In ZKBoost, the correctness of training is proved
        in a way similar to \textsc{Terrae},
        except that it relies on genric ZKPs for arithmetic circuits, which should be less efficient than our custom-designed proof system.
\end{enumerate}

Our implementation of ZKBoost strictly follows the above design choice of \cite{zkboost26}.
    We let our implementation of ZKBoost
    share the same cryptographic and 
    algebraic primitives as \textsc{Terrae},
    to ensure a fair comparison of the proof systems.

\paragraph{Implementation.} 
We implement \textsc{Terrae} and ZKBoost in C++.
  Cryptographic 
  operations are based on the mcl library~\cite{mcllib}, and the proof system is instantiated over the BLS12-381 pairing-friendly elliptic curve 
  with KZG polynomial commitments. This  
  instantiation targets roughly 128-bit security, following the standard security level of BLS12-381.
  We use Fiat-Shamir transformation to make 
  the proof systems non-interactive.

Our code is available at \url{https://anonymous.4open.science/r/ZK-GBDT-F601}.


\subsection{Asymptotic Comparison with Prior Works}

\Cref{tab:complexity-comparison} summarizes the asymptotic prover costs of
\textsc{Terrae} and the baseline. 

\begin{table}[t]
    \centering
    \small
    \caption{Asymptotic prover costs for tree-model training and inference
    certification. Here $h$ is the maximum root-to-leaf depth
    The
    $\tilde{O}(\cdot)$ notation hides logarithmic factors and fixed
    backend-dependent constants.
    Both verifiers run in polylogarithmic time.}
    \label{tab:complexity-comparison}
    \begin{tabular}{cccc}
        \toprule
        System & Proof of training & Proof of inference & Proof size\\
        \midrule
        Baseline
        & $\tilde{O}\bigl(TKMhd+N_{\mathrm{tot}}dB\bigr)$
        & $\tilde{O}(TKbhd)$
        & $O(1)$ \\
        \textsc{Terrae}
        & $\tilde{O}\bigl(TKMd+N_{\mathrm{tot}}dB\bigr)$
        & $\tilde{O}(TKbd)$
        & $O(1)$ \\
        \bottomrule
    \end{tabular}
\end{table}

One significant phenomenon is that the prover cost of training certification in \textsc{Terrae} is asymptotically smaller than the baseline by a factor of $h$.
This is due to the fact that \textsc{Terrae} avoids
committing the entire root-to-leaf path for proving inference, and thus avoids the need to prove the consistency of the path with the committed forest, which is a major cost in the baseline.

\subsection{Empirical Scaling in Three Dimensions}

We next study how proving time scales with three key workload parameters: the
number of iterations $T$, 
the maximum tree depth $d$, 
and the size of training dataset $N$.
We start with a basic case
$(T,d,N)=(512,1,512)$ and scale each parameter while keeping the others fixed.
For proof of training, we set the batch size $b=N$.
\Cref{tab:three-dim-scaling} reports prover time and peak prover memory for
both training and inference.

\begin{table*}[t]
    \centering
    % \scriptsize
    % \setlength{\tabcolsep}{2.5pt}
    \caption{Scaling results along three workload dimensions. Prover time is in
    seconds, peak memory is in GiB, and speedup is computed as
    $\text{ZKBoost}/\textsc{Terrae}$.}
    \label{tab:three-dim-scaling}
    \begin{tabular}{llrrrrrrrrrr}
        \toprule
        \multirow[c]{4}{*}{Scale} & \multirow[c]{4}{*}{Setting}
        & \multicolumn{5}{c}{Training}
        & \multicolumn{5}{c}{Inference} \\
        \cmidrule(lr){3-7} \cmidrule(lr){8-12}
        &
        & \multicolumn{3}{c}{Time (s)}
        & \multicolumn{2}{c}{Memory (GiB)}
        & \multicolumn{3}{c}{Time (s)}
        & \multicolumn{2}{c}{Memory (GiB)} \\
        \cmidrule(lr){3-5} \cmidrule(lr){6-7}
        \cmidrule(lr){8-10} \cmidrule(lr){11-12}
        &
        & \textsc{Terrae} & ZKBoost & Speedup
        & \textsc{Terrae} & ZKBoost
        & \textsc{Terrae} & ZKBoost & Speedup
        & \textsc{Terrae} & ZKBoost \\
        \midrule
        $T$-scale & $T=1$ & 28.93 & 127.76 & 4.42$\times$ & 1.79 & 12.30 & 22.02 & 121.93 & 5.54$\times$ & 1.50 & 4.20 \\
        $T$-scale & $T=2$ & 61.59 & 289.40 & 4.70$\times$ & 3.80 & 23.71 & 46.44 & 282.23 & 6.08$\times$ & 3.00 & 8.42 \\
        $T$-scale & $T=4$ & 136.03 & 634.10 & 4.66$\times$ & 8.37 & 47.29 & 104.20 & 585.24 & 5.62$\times$ & 5.86 & 16.87 \\
        $T$-scale & $T=8$ & 293.83 & 1354.88 & 4.61$\times$ & 19.75 & 91.47 & 223.96 & 1224.66 & 5.47$\times$ & 10.72 & 34.12 \\
        \midrule
        Depth-scale & $d=2$ & 32.75 & 293.56 & 8.96$\times$ & 1.98 & 22.29 & 22.44 & 280.64 & 12.50$\times$ & 1.51 & 8.33 \\
        Depth-scale & $d=3$ & 41.20 & 632.85 & 15.36$\times$ & 2.67 & 41.03 & 22.64 & 592.73 & 26.18$\times$ & 1.52 & 16.56 \\
        Depth-scale & $d=4$ & 63.29 & 1285.72 & 20.32$\times$ & 4.48 & 70.03 & 23.39 & 593.84 & 25.39$\times$ & 1.55 & 16.67 \\
        Depth-scale & $d=5$ & 134.90 & 1263.27 & 9.36$\times$ & 13.20 & 70.40 & 24.39 & 585.56 & 24.01$\times$ & 1.59 & 17.16 \\
        \midrule
        $N$-scale & $N=1024$ & 56.51 & 286.53 & 5.07$\times$ & 4.48 & 21.02 & 45.40 & 285.07 & 6.28$\times$ & 3.00 & 8.47 \\
        $N$-scale & $N=2048$ & 124.10 & 617.14 & 4.97$\times$ & 12.82 & 40.16 & 103.54 & 590.13 & 5.70$\times$ & 5.86 & 17.02 \\
        $N$-scale & $N=4096$ & 304.91 & 1226.27 & 4.02$\times$ & 41.32 & 76.96 & 218.96 & 1236.84 & 5.65$\times$ & 10.72 & 34.09 \\
        \bottomrule
    \end{tabular}
\end{table*}

\paragraph{Scaling with the number of trees.}
As expected, prover time in both systems grows roughly linearly with $T$, since
each additional tree introduces another proof instance of essentially the same
shape.  Across the entire sweep, \textsc{Terrae} consistently outperforms
ZKBoost by a stable margin: $4.42\times$--$4.70\times$ for training and
$5.47\times$--$6.08\times$ for inference.  From $T=1$ to $T=8$,
\textsc{Terrae}'s proving time increases from $28.93$\,s to $293.83$\,s for
training and from $22.02$\,s to $223.96$\,s for inference, which is close to
linear scaling.  Peak memory grows with $T$ as well, but \textsc{Terrae}
remains consistently lighter: at $T=8$, it uses $19.75$\,GiB for training and
$10.72$\,GiB for inference, compared with $91.47$\,GiB and $34.12$\,GiB for
ZKBoost.  This suggests that the improvement of \textsc{Terrae} is not tied to
a narrow parameter regime, but persists as the ensemble grows.

\paragraph{Scaling with tree depth.}
Depth is the dimension where the difference between the two approaches becomes
most pronounced.  On inference, \textsc{Terrae} is nearly insensitive to depth
in this range: its proving time changes only from $22.44$\,s at $d=2$ to
$24.39$\,s at $d=5$.  In contrast, ZKBoost remains around
$5.8\times 10^2$--$5.9\times 10^2$\,s, leading to speedups of
$12.50\times$--$26.18\times$.  A similar, though less uniform, trend appears in
training, where \textsc{Terrae} achieves $8.96\times$--$20.32\times$ speedups.
This behavior is consistent with the complexity discussion in
\Cref{tab:complexity-comparison}: ZKBoost pays heavily for depth-sensitive
route/path-consistency checks, whereas \textsc{Terrae} avoids proving explicit
root-to-leaf path objects and instead works with aggregated algebraic
constraints.  As a result, increasing depth has a much weaker effect on
\textsc{Terrae}, especially for inference.  The same effect appears in memory:
for inference, \textsc{Terrae} stays near $1.5$--$1.6$\,GiB throughout the depth
sweep, while ZKBoost rises to roughly $16$--$17$\,GiB once $d\geq 3$.

\paragraph{Scaling with the number of samples.}
When varying $N$, both systems again scale approximately linearly, reflecting
the fact that sample-dependent aggregation remains a dominant cost in both
designs.  Still, \textsc{Terrae} preserves a clear advantage throughout the
sweep: $4.02\times$--$5.07\times$ for training and $5.65\times$--$6.28\times$
for inference.  At the largest tested point, $N=4096$, \textsc{Terrae} reduces
training proving time from $1226.27$\,s to $304.91$\,s and inference proving
time from $1236.84$\,s to $218.96$\,s.  This indicates that the gains of
\textsc{Terrae} are maintained even when the workload is dominated by
sample-level processing.  Memory usage also increases with $N$, but the gap
remains visible at the largest tested point: for $N=4096$, \textsc{Terrae} uses
$41.32$\,GiB for training and $10.72$\,GiB for inference, while ZKBoost uses
$76.96$\,GiB and $34.09$\,GiB.

Overall, \Cref{tab:three-dim-scaling} shows a consistent pattern: \textsc{Terrae}
tracks the natural linear growth in $T$ and $N$, while substantially reducing
the constant factors, and it is markedly more robust to increasing depth.  The
depth sweep is particularly important in practice, since deeper trees are
exactly the setting in which path-based certification becomes expensive.  For
both systems, proof size remains essentially constant across these experiments,
and verifier time stays small; the main practical differences are therefore
prover time and peak prover memory, both of which favor \textsc{Terrae}.
